\documentclass[11pt]{article}

\usepackage[margin=1in]{geometry}
\usepackage{amsmath,amssymb}
\usepackage{graphicx}
\usepackage[numbers,sort]{natbib}
\usepackage[hidelinks]{hyperref}
\usepackage{booktabs}
\usepackage{multirow}
\usepackage{array}

\newcolumntype{L}[1]{>{\raggedright\arraybackslash}p{#1}}

\title{ResNet-18 Baseline Reproductions on Six Public\\Breast-Imaging Datasets: A Results Comparison}
\author{Nouran Fadlallah}
\date{}

\begin{document}
\maketitle

\begin{abstract}
We report single-modality binary (benign vs.\ malignant) classification
baselines on six public breast-imaging datasets: BUS-BRA
\citep{gomezflores2024busbra}, BUSI \citep{aldhabyani2020busi}, BUSC
\citep{busc2023mendeley,iqbalsharif2023pdfunet}, BrEaST-Lesions USG
\citep{pawlowska2024breast}, mini-MIAS \citep{suckling1994mias}, and
BreastDM \citep{zhao2023breastdm}. They span ultrasound, mammography, and
DCE-MRI. Every dataset uses a pretrained ResNet-18 under stratified group
$k$-fold cross validation, at patient or case level where the release gives
an identifier. For BUS-BRA we also reproduce the paper's own \textbf{ResNet-18 vs.\
ResNet-50} comparison, on its official fold assignment. Neither BUSI nor
BrEaST-Lesions USG reports its own classifier benchmark, so we add an
\textbf{EfficientNet-B0} comparison there instead, the backbone most used
in later work on these releases. We report accuracy, sensitivity,
specificity, precision, F1, and AUC for every run, and compare each result
to the source paper's own numbers where one exists.
\end{abstract}

\section{Datasets and Task}

All six datasets are reduced to one binary task: benign vs.\ malignant. We
drop modality-specific extras such as BI-RADS categories, normal-tissue
classes, and segmentation masks used only for ROI cropping. Table
\ref{tab:datasets} lists each dataset's size, modality, and split.

\begin{table}[htbp]
  \centering
  \caption{Dataset summary.}
  \label{tab:datasets}
  \resizebox{\textwidth}{!}{%
  \begin{tabular}{lL{3.4cm}rL{5cm}}
    \toprule
    Dataset & Modality & N & Split used \\
    \midrule
    BUS-BRA & Ultrasound & 1{,}875 imgs / 1{,}064 patients & Official case-level 5-fold (\texttt{K5B}) \\
    BUSI & Ultrasound & 647 imgs (437 benign / 210 malignant) & Our stratified 5-fold, row-level \\
    BUSC & Ultrasound & 250 imgs (100 / 150) & Our stratified 5-fold, row-level \\
    BrEaST-Lesions USG & Ultrasound & 256 cases & Our stratified 5-fold, case-level (\texttt{CaseID}) \\
    mini-MIAS & Mammography & 322 imgs, lesion patches & Our stratified 5-fold, patient-level \\
    BreastDM & DCE-MRI (\texttt{img9Se}) & 1{,}319 arrays & Official train/val split \\
    \bottomrule
  \end{tabular}%
  }
\end{table}

\section{Results}

\subsection{Training protocol}

Every backbone is ImageNet-pretrained. We train with SGD (momentum $0.9$,
learning rate $10^{-3}$) and class-weighted cross-entropy, for a fixed 100
epochs with no early stopping. We keep the checkpoint with the lowest
validation loss. We compute metrics on a held-out test fold that never
appears in training or checkpoint selection. Reported values are mean $\pm$
standard deviation across five folds. BreastDM is the exception: it uses
the dataset's own official train/validation split, with no cross-validation.
EfficientNet-B0 runs use the same optimizer, loss, and epoch budget as the
ResNet runs. All inputs are $224\times224$ except BUSC, which is
$128\times128$ to match the PDF-UNet release's stated resolution.

\subsection{Preprocessing}
\label{sec:preprocessing}

Preprocessing depends on what annotation each release provides (Table
\ref{tab:preprocessing}). BUS-BRA, BUSI, and BrEaST-Lesions USG each ship a
per-image tumor mask, so they share one pipeline: grayscale, $3\times3$
median filter, mask-guided ROI crop with a 5\% margin, train-time flips and
rotation ($\pm10^\circ$) and scale-crop, resize, an optional 5th/95th-
percentile contrast stretch, then replicate to 3 channels. We apply the
contrast stretch here, matching \citet{gomezflores2024busbra}. BUSC and
mini-MIAS skip it; their transform does not implement it. mini-MIAS instead
crops a lesion-centered patch at $1.5\times$ the annotated radius, using
the coordinates in \texttt{Info.txt}. BUSC's public release has no mask and
no lesion coordinates, so we resize the full image with no crop.
BreastDM's \texttt{img9Se} arrays are already derived by the dataset's own
pipeline. We only rescale them to $[0,1]$ and bilinearly resize each of the
nine slice channels. Augmentation there is flip-only, with no rotation, to
avoid interpolating across slice boundaries.

\begin{table}[htbp]
  \centering
  \caption{Preprocessing pipeline by dataset.}
  \label{tab:preprocessing}
  \resizebox{\textwidth}{!}{%
  \begin{tabular}{lL{10cm}}
    \toprule
    Dataset & Preprocessing \\
    \midrule
    BUS-BRA, BUSI, BrEaST-Lesions USG &
      Grayscale $\to$ $3\times3$ median filter $\to$ mask-guided ROI crop (5\% margin)
      $\to$ resize $\to$ 5th/95th-percentile contrast stretch $\to$ 3-channel replicate.
      Train-time: h/v flip, rotation $\pm10^\circ$, scale-crop $1.0$--$1.1\times$. \\
    BUSC &
      Grayscale $\to$ resize (no ROI crop: no mask or lesion coordinates in the release)
      $\to$ 3-channel replicate. Train-time: h flip, rotation $\pm10^\circ$. \\
    mini-MIAS &
      Grayscale $\to$ lesion-centered crop at $1.5\times$ the annotated radius
      (from \texttt{Info.txt} coordinates) $\to$ resize $\to$ 3-channel replicate.
      Train-time: h flip, rotation $\pm10^\circ$. \\
    BreastDM &
      \texttt{img9Se} array (already dataset-derived) $\to$ rescale to $[0,1]$
      $\to$ bilinear resize, 9 channels preserved. Train-time: h/v flip only. \\
    \bottomrule
  \end{tabular}%
  }
\end{table}

\subsection{Backbone comparison against the source paper}

BUS-BRA is the only one of the six source papers with its own second-
backbone comparison: ResNet-18 vs.\ ResNet-50, both on its own 5-fold split
(Table 3 of \citealp{gomezflores2024busbra}). Table \ref{tab:busbra}
reproduces that comparison on the same official folds.

\begin{table}[htbp]
  \centering
  \caption{BUS-BRA: ResNet-18 vs.\ ResNet-50, our reproduction vs.\ the source paper \citep{gomezflores2024busbra}, official 5-fold CV (mean $\pm$ std).}
  \label{tab:busbra}
  \resizebox{\textwidth}{!}{%
  \begin{tabular}{llcccccc}
    \toprule
    Backbone & Source & Acc. & Sens. & Spec. & Prec. & F1 & AUC \\
    \midrule
    \multirow{2}{*}{ResNet-18}
      & Paper & $0.865 \pm 0.020$ & $0.859 \pm 0.075$ & $0.868 \pm 0.014$ & --- & --- & $0.931 \pm 0.025$ \\
      & Ours  & $0.828 \pm 0.023$ & $0.822 \pm 0.044$ & $0.832 \pm 0.031$ & $0.702 \pm 0.034$ & $0.756 \pm 0.026$ & $0.903 \pm 0.018$ \\
    \multirow{2}{*}{ResNet-50}
      & Paper & $0.859$          & $0.808$          & $0.883$          & --- & --- & $0.920$ \\
      & Ours  & $0.828 \pm 0.021$ & $0.809 \pm 0.088$ & $0.840 \pm 0.055$ & $0.713 \pm 0.066$ & $0.751 \pm 0.032$ & $0.912 \pm 0.021$ \\
    \midrule
    \multicolumn{2}{l}{Senior sonographer (paper, Table 4)} & $0.827$ & $0.932$ & $0.776$ & --- & --- & --- \\
    \bottomrule
  \end{tabular}%
  }
\end{table}

The paper finds \textbf{ResNet-18 slightly ahead of ResNet-50}: higher
sensitivity and accuracy, lower specificity. It attributes this to
ResNet-50's larger capacity overfitting a small, single-institution set.
Our reproduction shows the \textbf{same sensitivity ordering}, ResNet-18 at
\textbf{0.822} vs.\ ResNet-50 at \textbf{0.809}, though our gap is smaller.
Both backbones land at the same accuracy, \textbf{0.828}. AUC is slightly
higher for ResNet-50 in both the paper and our run.

\subsection{Backbones without a paper-reported comparison}

BUSI and BrEaST-Lesions USG are dataset-release papers. Neither reports its
own classifier benchmark. Table \ref{tab:effnet} adds an EfficientNet-B0
comparison for both instead, matching what later work on these releases
typically evaluates. This comparison is literature-informed, not
paper-reported.

\begin{table}[htbp]
  \centering
  \caption{BUSI and BrEaST-Lesions USG: ResNet-18 vs.\ EfficientNet-B0, our 5-fold CV (mean $\pm$ std). Neither source paper reports its own classifier benchmark.}
  \label{tab:effnet}
  \resizebox{\textwidth}{!}{%
  \begin{tabular}{llcccccc}
    \toprule
    Dataset & Backbone & Acc. & Sens. & Spec. & Prec. & F1 & AUC \\
    \midrule
    \multirow{2}{*}{BUSI \citep{aldhabyani2020busi}}
      & ResNet-18       & $0.884 \pm 0.038$ & $0.833 \pm 0.066$ & $0.908 \pm 0.026$ & $0.814 \pm 0.055$ & $0.823 \pm 0.059$ & $0.945 \pm 0.017$ \\
      & EfficientNet-B0 & $0.881 \pm 0.022$ & $0.871 \pm 0.032$ & $0.886 \pm 0.024$ & $0.787 \pm 0.039$ & $0.826 \pm 0.031$ & $0.938 \pm 0.016$ \\
    \multirow{2}{*}{BrEaST-Lesions USG \citep{pawlowska2024breast}}
      & ResNet-18       & $0.726 \pm 0.042$ & $0.728 \pm 0.201$ & $0.726 \pm 0.139$ & $0.654 \pm 0.101$ & $0.664 \pm 0.075$ & $0.836 \pm 0.079$ \\
      & EfficientNet-B0 & $0.694 \pm 0.096$ & $0.757 \pm 0.128$ & $0.655 \pm 0.152$ & $0.604 \pm 0.127$ & $0.661 \pm 0.091$ & $0.781 \pm 0.108$ \\
    \bottomrule
  \end{tabular}%
  }
\end{table}

On BUSI, the two backbones are \textbf{within noise of each other} on every
metric. ResNet-18 has higher specificity and AUC; EfficientNet-B0 has higher
sensitivity and F1. Neither is a clear winner at this dataset size. On
BrEaST-Lesions USG, ResNet-18 leads on accuracy, specificity, and AUC;
EfficientNet-B0 has slightly higher sensitivity. Both backbones carry
\textbf{high fold-to-fold variance} (ResNet-18 sensitivity std
$=\mathbf{0.201}$; EfficientNet-B0 accuracy std $=0.096$, AUC std $=0.108$),
from only about 51 test cases per fold. Read this ordering as
\textbf{directional, not a reliable ranking} (Section~\ref{sec:limitations}).

\subsection{Remaining datasets: ResNet-18 only}

BUSC, mini-MIAS, and BreastDM use ResNet-18 only. BUSC's own paper
\citep{iqbalsharif2023pdfunet} is a segmentation method with no
classification benchmark. mini-MIAS predates deep learning
\citep{suckling1994mias}; its paper reports no CNN result. BreastDM's paper
proposes a custom fusion architecture, LG-CAFN, not a swappable backbone.
Its 9-channel MRI input also only works with this codebase's ResNet family,
since non-ResNet backbones skip first-layer channel adaptation. Table
\ref{tab:resnet-only} reports these three against whatever comparison each
source provides.

\begin{table}[htbp]
  \centering
  \caption{BUSC, mini-MIAS, BreastDM: ResNet-18 only, vs.\ the source literature where a comparable number exists.}
  \label{tab:resnet-only}
  \resizebox{\textwidth}{!}{%
  \begin{tabular}{llcccccc}
    \toprule
    Dataset & Source & Acc. & Sens. & Spec. & Prec. & F1 & AUC \\
    \midrule
    \multirow{2}{*}{BUSC \citep{busc2023mendeley}}
      & Ours (5-fold) & $0.992 \pm 0.010$ & $1.000 \pm 0.000$ & $0.980 \pm 0.025$ & $0.987 \pm 0.016$ & $0.993 \pm 0.008$ & $0.9997 \pm 0.0007$ \\
      & Paper & \multicolumn{6}{l}{No classification benchmark (segmentation-only paper \citep{iqbalsharif2023pdfunet})} \\
    \multirow{2}{*}{mini-MIAS \citep{suckling1994mias}}
      & Ours (5-fold, patient-level) & $0.599 \pm 0.096$ & $0.496 \pm 0.233$ & $0.675 \pm 0.098$ & $0.496 \pm 0.161$ & $0.486 \pm 0.205$ & $0.597 \pm 0.107$ \\
      & Paper & \multicolumn{6}{l}{No CNN benchmark (1994, predates deep learning)} \\
    \multirow{2}{*}{BreastDM \citep{zhao2023breastdm}}
      & Ours (official split) & $0.880$ & $0.968$ & $0.542$ & $0.891$ & $0.928$ & $0.907$ \\
      & Paper (LG-CAFN) & $0.882$ & --- & --- & --- & --- & $0.915$ \\
    \bottomrule
  \end{tabular}%
  }
\end{table}

BUSC's \textbf{near-perfect result}, across every metric and every fold, is
\textbf{a flag, not a win}. Cross-validation rules out a single-split
fluke. But 250 images with no patient identifier, at $128\times128$
resolution, leave open whether benign and malignant images are trivially
separable by something other than the lesion (Section~\ref{sec:limitations}).
mini-MIAS's cross-validated result is \textbf{the honest number} for this
task framing. A single earlier evaluation, on one ad-hoc split, showed
90\% accuracy; cross-validation shows that was \textbf{a lucky split, not a
stable result}. BreastDM's plain ResNet-18 lands within \textbf{0.002
accuracy and 0.008 AUC} of the paper's own fusion architecture, LG-CAFN.
Its specificity is much weaker, \textbf{0.542}: the model
\textbf{over-calls malignant}.

\section{Discussion: Is the Preprocessing Sufficient?}
\label{sec:preprocessing-discussion}

Section~\ref{sec:preprocessing} described two pipelines: a mask-guided ROI
crop with contrast stretching (BUS-BRA, BUSI, BrEaST-Lesions USG), and a
lighter pipeline with no crop or contrast stretch (BUSC, mini-MIAS, and
BreastDM's already-derived arrays). Tables~\ref{tab:busbra}--\ref{tab:resnet-only}
let us ask, per dataset, whether that choice was adequate or needs further
investigation.

\paragraph{Where preprocessing looks adequate.} BUS-BRA reproduces the
paper's own 5-fold numbers within about \textbf{0.03--0.04 accuracy and
AUC} (Table~\ref{tab:busbra}), using the mask-guided crop and contrast
stretch. We can already explain the remaining gap: our pipeline replicates
grayscale into 3 channels, while the paper encodes the mask into R/G/B
channels (Section~\ref{sec:limitations}). \textbf{The crop and
contrast-stretch steps are not the cause.} BUSI reaches the literature's
commonly cited \textbf{90--96\%} accuracy range with the same pipeline,
unchanged. BreastDM needs almost no preprocessing of its own, since the
release's \texttt{img9Se} arrays are already derived, and still lands
within \textbf{0.002 accuracy} of the paper's purpose-built fusion
architecture. In these three cases, more preprocessing effort means results
closer to the literature. \textbf{The numbers give no reason for further
preprocessing work here.}

\paragraph{Where preprocessing is a real suspect.} BUSC is the opposite
signal. It has the \textbf{least preprocessing} of any dataset here: no ROI
crop at all, since the release ships neither a mask nor lesion coordinates
(Table~\ref{tab:preprocessing}). It also has the \textbf{best result by a
wide margin}: \textbf{0.992--0.9997} across every metric and every fold
(Table~\ref{tab:resnet-only}). This is the pattern we would expect if the
model is keying on something other than the lesion, since there is no crop
step to exclude background context. A device-specific border, an
annotation artifact, or an acquisition setting could correlate with the
class label in this 250-image release. \textbf{More cross-validation will
not resolve this.} It needs a \textbf{Grad-CAM or saliency check} on the
trained model. If a lesion boundary can be recovered by other means, such
as manual annotation or a generic saliency-based ROI heuristic, a re-run
with an actual crop would show whether accuracy drops toward the other
ultrasound datasets' range. Until that check happens, \textbf{treat BUSC's
number as unverified}, not as a working classifier.

\paragraph{Where the picture is genuinely unclear.} mini-MIAS and
BrEaST-Lesions USG both crop before classification: a radius-based lesion
patch for mini-MIAS, a mask-guided ROI for BrEaST. Both still produce
\textbf{weak or unstable results}. mini-MIAS gets \textbf{$0.599 \pm 0.096$}
accuracy; BrEaST's sensitivity std reaches \textbf{0.201}. Both are also
the \textbf{two smallest datasets} by patient or case count, so we cannot
yet tell whether preprocessing (crop margin, patch size, missing contrast
stretch) or sample size causes the weak results. Two follow-ups would
separate these causes. First, apply the BUS-BRA-style contrast stretch to
mini-MIAS, currently skipped (Table \ref{tab:preprocessing}); this is a
one-line change in \texttt{MiasTransform}. Second, for mini-MIAS, compare
the current lesion-patch framing against a full-breast-image baseline.
\citet{suckling1994mias} does not require the patch framing, and
\texttt{dataset\_reproduction\_plan} lists it as an open design choice.
Until we run one of these, \textbf{we cannot rule preprocessing in or out
as the bottleneck here}.

\section{Mixed-Backbone Multimodal Fusion}
\label{sec:mixed-fusion}

All results so far are single-modality baselines. A separate follow-up
study, \texttt{docs/mixed\_backbone\_results.md}, trains the repository's
3-branch late-fusion model instead: \texttt{FusionLateModel}, with
mammography from mini-MIAS, ultrasound from BUS-BRA, BUSI, BUSC, and
BrEaST-Lesions USG, and MRI from BreastDM. It varies the ultrasound
branch's backbone and initialization to separate three effects an earlier
single run had confounded: ResNet-18 vs.\ ResNet-50, warm-starting from a
converged single-dataset checkpoint, and RadImageNet vs.\ ImageNet
pretraining. Mammography and MRI stay fixed across all three runs:
ResNet-18, warm-started from \texttt{runs/run-1}'s single-dataset
checkpoints. Only the ultrasound branch varies. All four runs share one
combined validation set of 427 samples (10 mammography, 300 ultrasound, 117
MRI) and use the default recipe, not \texttt{paper-match}: Adam, a 10-epoch
budget, early-stopping patience 5. This is not the SGD/weighted-CE protocol
used for the single-dataset results above.

\begin{table}[htbp]
  \centering
  \caption{Combined 3-branch fusion: varying only the ultrasound branch's backbone/initialization. Selection is by lowest validation loss, not by accuracy.}
  \label{tab:mixed-fusion}
  \resizebox{\textwidth}{!}{%
  \begin{tabular}{lllcccc}
    \toprule
    Run & Ultrasound backbone & Ultrasound init & Best epoch & Best val loss & Val acc. & Epochs run \\
    \midrule
    Baseline & ResNet-18 & ImageNet, trained from scratch & --- & $0.467$ & $0.741$ & 30 \\
    A & ResNet-50 & Warm-started (converged BUS-BRA fold-1 model) & 4 & $0.421$ & $0.808$ & 9 (early-stopped) \\
    B & ResNet-18 & Warm-started (converged BUS-BRA fold-1 model) & 10 & $\mathbf{0.395}$ & $\mathbf{0.836}$ & 10 (full budget) \\
    C & ResNet-50 & RadImageNet init, trained in-loop & 2 & $0.428$ & $0.815$ & 7 (early-stopped) \\
    \bottomrule
  \end{tabular}%
  }
\end{table}

Holding the two warm-started branches fixed and varying only the ultrasound
branch isolates three effects:

\begin{enumerate}
  \item \textbf{Warm-starting from a converged single-dataset model is the
    main effect.} Both warm-started runs (A: \textbf{0.421}, B:
    \textbf{0.395}) beat both non-warm-started runs (baseline:
    \textbf{0.467}, C: \textbf{0.428}) by a similar margin. Backbone and
    pretraining source do not change this.
  \item \textbf{RadImageNet beats ImageNet when training in-loop.}
    Comparing the two non-warm-started runs directly, baseline (ImageNet)
    vs.\ Run C (RadImageNet), both trained from scratch inside the fusion
    loop, RadImageNet initialization helps. This is the \textbf{clearest
    evidence so far} that RadImageNet adds value on its own.
  \item \textbf{ResNet-18 edges out ResNet-50 under the same warm-start
    treatment} (Run B \textbf{0.395} vs.\ Run A \textbf{0.421}). This
    matches the BUS-BRA paper's own finding (Table~\ref{tab:busbra}) and
    this report's single-mode BUS-BRA 5-fold results. \textbf{The
    ResNet-50 swap is not the source of improvement} over the baseline.
    Under matched warm-start treatment, it is a slight drag compared to
    ResNet-18.
\end{enumerate}

\paragraph{Caveats specific to this study.} The mammography branch's
per-epoch accuracy comes from only \textbf{10 validation samples}. Read it
as noise in every run. The ultrasound warm-start (Runs A/B) uses only
\textbf{one of the five} BUS-BRA cross-validation folds, fold 1, not an
aggregate. None of these four runs use the \texttt{paper-match} recipe from
Section~\ref{sec:preprocessing-discussion} onward, so they are \textbf{not
as rigorous} as this report's single-dataset numbers. Early stopping also
cuts the three runs at different points: 9, 10, and 7 epochs. Run B hit the
full 10-epoch budget still improving, so its true optimum may be higher.
\textbf{This is a fusion-mechanism ablation, not a standalone benchmark
result.}

\section{Limitations}
\label{sec:limitations}

\begin{itemize}
  \item \textbf{Split granularity.} BUS-BRA, BrEaST-Lesions USG, and
    mini-MIAS use patient- or case-level folds, via \texttt{K5B},
    \texttt{CaseID}, and \texttt{patient}. BUSI and BUSC have no patient
    identifier in their public release, so we split them at the image
    level. A fold boundary can still place near-duplicate crops of the
    same lesion on both sides.
  \item \textbf{BUS-BRA input encoding.} Our pipeline repeats the
    grayscale, contrast-stretched, mask-cropped image into three identical
    channels. The source paper instead encodes the lesion mask into
    separate R/G/B channels. This likely explains part of the remaining
    \textbf{0.03--0.04} accuracy/AUC gap in Table~\ref{tab:busbra}.
  \item \textbf{Single seed per fold.} All reported std is across the 5
    folds, not across repeated seeds within a fold. The paper's own
    reported std for BUS-BRA is also a single-seed, 5-fold figure.
  \item \textbf{Dataset scale.} mini-MIAS (about 130 lesion patches total)
    and BrEaST-Lesions USG (256 cases) are small enough that 5-fold CV
    still leaves high per-fold variance. Read results there as
    directional.
\end{itemize}

\bibliographystyle{plainnat}
\bibliography{references}

\end{document}
